% SPDX-License-Identifier: CC-BY-4.0
\documentclass[11pt,letterpaper]{article}

\usepackage[T1]{fontenc}
\usepackage{amsmath,amssymb}
\usepackage{tikz}
\usepackage{pgfplots}
\pgfplotsset{compat=1.18}
\usepackage{algorithm}
\usepackage{algpseudocode}
\usepackage{booktabs}
\usepackage{enumitem}
\usepackage[hidelinks]{hyperref}

\newcommand{\N}{\mathbb{N}}
\newcommand{\Z}{\mathbb{Z}}
\newcommand{\Q}{\mathbb{Q}}
\newcommand{\R}{\mathbb{R}}
\newcommand{\C}{\mathbb{C}}
\newcommand{\Map}[3]{#1 \colon #2 \to #3}
\newcommand{\Mod}[1]{#1\text{-}\mathrm{Mod}}
\newcommand{\GroupHom}{\mathrm{Hom}_{\mathbf{Grp}}}
\newcommand{\RingHom}{\mathrm{Hom}_{\mathbf{Ring}}}

\newtheorem{theorem}{Theorem}
\newtheorem{lemma}[theorem]{Lemma}
\newtheorem{definition}[theorem]{Definition}
\newtheorem{remark}[theorem]{Remark}
\newenvironment{proof}{\noindent\textit{Proof.}}{\hfill\(\square\)}

\title{Generalized Collatz Library}
\author{ascendnoosphere}
\date{\today}

\begin{document}
\maketitle

\begin{abstract}
Let $q,D\geq2$ and let the positive homogeneous raw map divide by $D$ when
possible and multiply by $q$ otherwise.  Suppose $q^a=D^b$ for some positive
$a,b$.  A primitive common-base normal form $q=t^u$, $D=t^v$ converts the
integer dynamics into an exact exponent recurrence.  Its transient has a
strictly decreasing exponent, and its recurrent part is a coprime rotation
on $u+v$ residues.  Consequently every positive orbit has an eventually
primitive-periodic tail of raw period $u+v$, and every raw iterate is at most
$qDn$.  The proof reconstructs every integer state at the original raw clock.
\end{abstract}

\section{Raw dynamics and the target}

\begin{definition}[Positive homogeneous raw map]
For $q,D\in\N$ with $q,D\geq2$, define
\[
 T_{q,D}(n)=
 \begin{cases}
  n/D,&D\mid n,\\
  qn,&D\nmid n
 \end{cases}
 \qquad(n\in\N_{>0}).
\]
The divisibility test is made before every instruction.  Thus one division is
one raw step, even if further divisions immediately follow.
\end{definition}

\begin{definition}[Multiplicative dependence]
The coefficients $q,D\geq2$ are multiplicatively dependent if there are
$a,b\in\N_{>0}$ such that
\[
 q^a=D^b.
\]
\end{definition}

The obstruction is the unbounded integer state.  The reduction below removes
that obstruction without changing the clock: an invariant factor is retained
exactly, while a single natural-number exponent first decreases and then
enters a finite rotation.

\section{Primitive common-base normal form}

\begin{lemma}[Common-base extraction]
\label{lem:normal-form}
If $q,D\geq2$ and $q^a=D^b$ for positive $a,b$, then there exist
$t,u,v\in\N$ such that
\[
 t\geq2,\qquad u,v\geq1,\qquad \gcd(u,v)=1,
 \qquad q=t^u,\qquad D=t^v.
\]
\end{lemma}

\begin{proof}
Let $g=\gcd(a,b)$ and write $a=gv$ and $b=gu$.  Then
\[
 (q^v)^g=(D^u)^g.
\]
Positive integer powers are injective, so $q^v=D^u$, and
$\gcd(u,v)=1$.

For each prime $p$, let $\alpha_p$ and $\beta_p$ be its exponents in $q$ and
$D$.  Unique prime factorization gives
\[
 v\alpha_p=u\beta_p.
\]
Since $u$ and $v$ are coprime, $u\mid\alpha_p$ and $v\mid\beta_p$.  Write
$\alpha_p=u\gamma_p$ and $\beta_p=v\gamma_p$.  The finite product
\[
 t=\prod_p p^{\gamma_p}
\]
then satisfies $q=t^u$ and $D=t^v$.  Because $q\geq2$ and $u>0$, one has
$t\geq2$.  This is the standard prime-factorization argument; see
\cite[Chapter I]{HardyWright08}.
\end{proof}

Dividing $a,b$ by their gcd is essential.  Without the coprimality condition,
the exponent rotation can return before the displayed nonprimitive length.

\section{Exact exponent dynamics}

Fix a primitive presentation
\[
 q=t^u,\qquad D=t^v,
 \qquad t\geq2,\quad u,v\geq1,\quad\gcd(u,v)=1.
\]

\begin{lemma}[Canonical factorization]
\label{lem:canonical-factorization}
For every $n\in\N_{>0}$ there is a unique pair $(r,k)$ with
\[
 n=rt^k,\qquad k\in\N,\qquad t\nmid r.
\]
\end{lemma}

\begin{proof}
The set $\{e\in\N:t^e\mid n\}$ is nonempty and finite, because $t\geq2$.
Let $k$ be its maximum and put $r=n/t^k$.  Maximality gives $t\nmid r$.
If $rt^k=r't^{k'}$ with both $r,r'$ not divisible by $t$, then $k<k'$
would imply $t\mid r$, and $k'<k$ would imply $t\mid r'$.  Hence $k=k'$
and then $r=r'$.
\end{proof}

\begin{definition}[Exponent map]
Define $F\colon\N\to\N$ by
\[
 F(k)=
 \begin{cases}
  k-v,&v\leq k,\\
  k+u,&k<v.
 \end{cases}
\]
\end{definition}

\begin{lemma}[One-step reconstruction]
\label{lem:one-step}
If $n=rt^k$ with $t\nmid r$, then
\[
 T_{q,D}(n)=rt^{F(k)}.
\]
In particular, the factor $r$ is invariant.
\end{lemma}

\begin{proof}
Because $D=t^v$ and $t\nmid r$, one has
\[
 D\mid rt^k\quad\Longleftrightarrow\quad v\leq k.
\]
When $v\leq k$, the division branch gives
$rt^k/t^v=rt^{k-v}$.  When $k<v$, the affine branch gives
$t^u rt^k=rt^{k+u}$.  These are exactly the two clauses defining $F$.
\end{proof}

\begin{theorem}[Exact raw reconstruction]
\label{thm:reconstruction}
Let $n=rt^{k_0}$ be the canonical factorization.  For every $j\in\N$,
\[
 T_{q,D}^{j}(n)=rt^{F^j(k_0)}.
\]
The exponent on the right is evaluated after exactly $j$ raw instructions.
\end{theorem}

\begin{proof}
The assertion at $j=0$ is the canonical factorization.  If it holds at $j$,
Lemma~\ref{lem:one-step} applied to $rt^{F^j(k_0)}$ gives
\[
 T_{q,D}^{j+1}(n)=rt^{F(F^j(k_0))}=rt^{F^{j+1}(k_0)}.
\]
Induction proves the formula for every $j$.
\end{proof}

\section{Finite recurrent rotation}

Put $m=u+v$.

\begin{lemma}[Entry into the recurrent interval]
\label{lem:entry}
For every $k\in\N$, some iterate $F^s(k)$ lies in
$\{0,1,\ldots,m-1\}$.
\end{lemma}

\begin{proof}
If $k<m$, take $s=0$.  If $k\geq m$, then $k\geq v$ and
$F(k)=k-v<k$.  Strong induction on $k$ therefore supplies a finite entry
time.  The decreasing resource is the exponent itself, and each unresolved
noncore step lowers it by the fixed positive amount $v$.
\end{proof}

\begin{lemma}[Rotation formula]
\label{lem:rotation}
If $0\leq k<m$, then
\[
 F(k)\equiv k+u\pmod m
 \qquad\text{and}\qquad 0\leq F(k)<m.
\]
Consequently, for every $j\in\N$,
\[
 F^j(k)\equiv k+ju\pmod m.
\]
\end{lemma}

\begin{proof}
If $k<v$, then $F(k)=k+u<m$.  If $v\leq k<m$, then
\[
 F(k)=k-v=k+u-(u+v)=k+u-m.
\]
This proves the one-step congruence and interval preservation.  The iterated
formula follows by induction on $j$.
\end{proof}

\begin{lemma}[Primitive exponent period]
\label{lem:primitive-exponent}
For every $0\leq k<m$,
\[
 F^m(k)=k.
\]
If $0<j<m$, then $F^j(k)\neq k$.
\end{lemma}

\begin{proof}
Lemma~\ref{lem:rotation} gives
$F^m(k)\equiv k+mu\equiv k\pmod m$.  Both terms lie in the recurrent
interval, so they are equal.  Conversely, if $F^j(k)=k$, then
$m\mid ju$.  Since
\[
 \gcd(u,m)=\gcd(u,u+v)=\gcd(u,v)=1,
\]
one has $m\mid j$.  This is impossible when $0<j<m$.
\end{proof}

\begin{theorem}[Eventual primitive periodicity]
\label{thm:periodicity}
Every positive orbit of $T_{q,D}$ enters a raw cycle of primitive period
\[
 u+v.
\]
More precisely, for every $n>0$ there is an $s\in\N$ such that
\[
 T_{q,D}^{u+v}(T_{q,D}^s(n))=T_{q,D}^s(n),
\]
and no positive $j<u+v$ returns the tail state.
\end{theorem}

\begin{proof}
Write $n=rt^{k_0}$.  By Lemma~\ref{lem:entry}, choose $s$ so that
$k_s=F^s(k_0)<u+v$.  Theorem~\ref{thm:reconstruction} and
Lemma~\ref{lem:primitive-exponent} give the return after $u+v$ raw steps.
If an earlier positive return occurred, cancellation of the positive factor
$r$ would give $t^{F^j(k_s)}=t^{k_s}$.  Since $t\geq2$, equality of powers
forces $F^j(k_s)=k_s$, contradicting primitive exponent period.
\end{proof}

\section{Uniform height}

\begin{lemma}[Invariant exponent bound]
\label{lem:exponent-bound}
For all $k,j\in\N$,
\[
 F^j(k)<\max\{k+1,u+v\}.
\]
\end{lemma}

\begin{proof}
Let $B=\max\{k+1,u+v\}$.  Initially $k<B$.  Suppose $x<B$.
If $v\leq x$, then $F(x)=x-v<x<B$.  If $x<v$, then
$F(x)=x+u<u+v\leq B$.  Thus the strict bound is invariant, and induction on
$j$ proves the claim.
\end{proof}

\begin{theorem}[Source-relative height bound]
\label{thm:height}
For every $n>0$ and every raw time $j$,
\[
 T_{q,D}^j(n)\leq qDn.
\]
\end{theorem}

\begin{proof}
Write $n=rt^{k_0}$.  Lemma~\ref{lem:exponent-bound} implies
$F^j(k_0)\leq k_0+u+v$.  Hence
\[
 T_{q,D}^j(n)
 =rt^{F^j(k_0)}
 \leq rt^{k_0+u+v}
 =t^{u+v}n
 =qDn.
\]
This includes the transient and the periodic tail.
\end{proof}

\section{Reduction certificate and scope}

The original obligation is the orbit classification for every positive seed
under zero offset, $q,D\geq2$, and a positive power relation.  The replacement
obligations are common-base extraction, exponent recurrence, and exact
reconstruction.  Lemma~\ref{lem:entry} supplies a well-founded decrease, and
Lemmas~\ref{lem:rotation}--\ref{lem:primitive-exponent} solve the finite
recurrent problem.  Theorem~\ref{thm:reconstruction} is the recovery map and
retains every raw instruction.  Thus no convergence argument, accelerated
clock, or unproved all-seed assertion is used.

The Lean development proves the common-base normal form, the modular exponent
formula, entry into the recurrent interval, exact state reconstruction,
primitive tail period, and the height bound.  It also identifies the natural
map with the positive restriction of the signed foundational raw map.  The Go
checker constructs the transient and all $u+v$ cycle states with
arbitrary-precision integers; its bounded executable family checks $61{,}250$
certificates and a separate $42$-step large-integer cycle.

The theorem excludes nonzero offsets, zero and negative seeds, signed
moduli, and coefficients below two.  It does not imply universal classical
Collatz capture.  Paired with the multiplicatively independent escape theorem,
it completes the positive homogeneous dichotomy for $q,D\geq2$.

\bibliographystyle{amsalpha}
\bibliography{DependentHomogeneous}

\end{document}
